\medskip

Voici la répartition des principaux groupes sanguins des habitants de France:

\begin{center}
\psset{xunit=1.2cm,yunit=0.25cm}
\begin{pspicture}(8.5,41)
\psaxes[linewidth=1.25pt,Dx=10,Dy=100](0,0)(0,0)(8.5,40)
\uput[d](1,0){A+}\uput[d](2,0){O+}\uput[d](3,0){B+}\uput[d](4,0){A$-$}
\uput[d](5,0){O$-$}\uput[d](6,0){AB+}\uput[d](7,0){B$-$}\uput[d](8,0){AB$-$}
\multido{\n=0+5}{9}{\uput[l](0,\n){\n\,\%}}
\psframe[fillstyle=solid,fillcolor=lightgray](0.8,0)(1.2,38.2)\uput[u](1,38.2){38,2\,\%}
\psframe[fillstyle=solid,fillcolor=lightgray](1.8,0)(2.2,36.5)\uput[u](2,36.5){36,5\,\%}
\psframe[fillstyle=solid,fillcolor=lightgray](2.8,0)(3.2,7.7)\uput[u](3,7.7){7,7\,\%}
\psframe[fillstyle=solid,fillcolor=lightgray](3.8,0)(4.2,6.8)\uput[u](4,6.8){6,8\,\%}
\psframe[fillstyle=solid,fillcolor=lightgray](4.8,0)(5.2,6.5)\uput[u](5,6.5){6,5\,\%}
\psframe[fillstyle=solid,fillcolor=lightgray](5.8,0)(6.2,2.5)\uput[u](6,2.5){2,5\,\%}
\psframe[fillstyle=solid,fillcolor=lightgray](6.8,0)(7.2,1.4)\uput[u](7,1.4){1,4\,\%}
\psframe[fillstyle=solid,fillcolor=lightgray](7.8,0)(8.2,0.4)\uput[u](8,0.4){0,4\,\%}
\end{pspicture}

\medskip

\emph{\scriptsize Source : https://fr.statista.com!statistiques/6 56036/ groupes-sanguins-repartition-rh- france/}

\end{center}

A+, O+, B+, A$-$, O$-$, AB+, B$-$ et AB$-$ sont les différents groupes sanguins combinés aux rhésus.

Par exemple: A + est le groupe sanguin A de rhésus +.

Une expérience aléatoire consiste à choisir une personne au hasard dans la population française et à déterminer son groupe sanguin et son rhésus.

Dans l'exercice, on adopte les notations du type :

$A$ + est l'évènement \og la personne est de groupe sanguin A et de rhésus + \fg 

$A -$ est l'évènement \og la personne est de groupe sanguin A et de rhésus $-$ \fg

$A$ est l'évènement \og la personne est de groupe sanguin A \fg

\medskip

Les parties 1 et 2 sont indépendantes.


\textbf{Partie 1}

\medskip

On note $Rh$ + l'évènement \og La personne est de rhésus positif \fg.

\begin{enumerate}
\item Justifier que la probabilité que la personne choisie soit de rhésus positif est égale à 0,849.
\item Démontrer à l'aide des données de l'énoncé que $P_{Rh+}(A) = 0,450$ à 0,001 près.
\item Une personne se souvient que son groupe sanguin est AB mais a oublié son rhésus.

Quelle est la probabilité que son rhésus soit négatif ? Arrondir le résultat à 0,001 près.
\end{enumerate}

\bigskip

\textbf{Partie 2}

\medskip

Dans cette partie, les résultats seront arrondis à $0,001$ près.

\medskip

Un donneur universel de sang est une personne de groupe sanguin O et de rhésus négatif. On rappelle que 6,5\,\% de la population française est de groupe O$-$.

\medskip

\begin{enumerate}
\item On considère $50$ personnes choisies au hasard dans la population française et on note $X$ la variable aléatoire qui compte le nombre de donneurs universels.
	\begin{enumerate}
		\item Déterminer la probabilité que 8 personnes soient des donneurs universels. Justifier votre réponse.
		\item On considère la fonction ci-dessous nommée \texttt{proba} d'argument k écrite en langage Python.
		
\begin{center}
\begin{tabular}{|l|}\hline
def proba(k):\\
\quad p=0\\
\quad for i in range(k+1):\\
\quad \quad p = p + binomiale(i,50,0.065)\\ 
\quad return p\\ \hline
\end{tabular}
\end{center}
Cette fonction utilise la fonction binomiale d'argument $i, n$ et $p$, créée pour l'occasion, qui renvoie la valeur de la probabilité $P(X = i)$ dans le cas où $X$ suit une loi binomiale de paramètres $n$ et $p$.

Déterminer la valeur numérique renvoyée par la fonction \texttt{proba} lorsqu'on saisit \texttt{proba(8)} dans la console Python. Interpréter ce résultat dans le contexte de l'exercice.
	\end{enumerate}
\item Quel est le nombre minimal de personnes à choisir au hasard dans la population française pour que la probabilité qu'au moins une des personnes choisies soit donneur universel, soit supérieure à $0,999$.
\end{enumerate}

\bigskip
